\subsection*{Objetivos Específicos}
\addcontentsline{toc}{subsection}{Objetivo Específicos}

\begin{itemize}
	\item Realizar un diagnóstico del contexto macroeconómico y del comportamiento financiero en Bogotá para identificar las necesidades críticas de alfabetización y gestión de capital en segmentos como jóvenes (18-28 años), trabajadores informales y nómadas digitales.
	
	\item Estructurar el modelo de negocio mediante la metodología CANVAS, definiendo una propuesta de valor basada en la educación financiera profunda y diseñando estrategias de monetización \textit{freemium} que garanticen la sostenibilidad operativa del emprendimiento.
	
	\item Definir los requerimientos funcionales y no funcionales del sistema, estableciendo una arquitectura técnica modular basada en microservicios que soporte la escalabilidad horizontal y la integración de modelos de Machine Learning para rutas de aprendizaje adaptativas.
	
	\item Elaborar un análisis financiero proyectado a cinco años que permita evaluar la viabilidad económica del proyecto, determinando el punto de equilibrio, los indicadores de rentabilidad y la estructura de costos asociada a la infraestructura cloud y el talento humano.
	
	\item Desarrollar un Producto Mínimo Viable (PMV) que integre las funcionalidades de simulación financiera y gestión de flujo de caja, validando la calidad del software y la experiencia de usuario (UX) mediante una metodología de desarrollo ágil que permita iteraciones basadas en el \textit{feedback} de los usuarios.
	
	\item Formular el análisis de impactos sociales, económicos, ambientales y tecnológicos del proyecto, estableciendo indicadores de seguimiento proyectados a corto, mediano y largo plazo para asegurar la alineación del servicio con los objetivos de desarrollo sostenible y el ecosistema Open Finance.
\end{itemize}